\chapter{Convex Structure}\label{ch:convex_structure}

This chapter follows the convex-optimization framework of
\cite[Chapter~4]{Wolf2012Quantum}.

\begin{definition}[Primal and dual conic programs]
    \label{def:conic_programs}
    \lean{ConicProgram.primalFeasible, ConicProgram.dualFeasible,
        ConicProgram.primalValue, ConicProgram.dualValue}
    \leanok
    Let $V$ and $V'$ be finite-dimensional real Hilbert spaces, let
    $K\subseteq V$ be a closed pointed convex cone with nonempty interior, let
    $T:V\to V'$ be linear, and fix $c\in V$ and $b\in V'$.  As in
    \cite[Chapter~4, equations~(4.1)--(4.2)]{Wolf2012Quantum}, define the dual
    cone, the primal and dual feasible sets, and their values by
    \begin{align}
        K^*
         & :=\{u\in V:\forall z\in K,\ \langle z|u\rangle\geq0\},
        \label{eq:conic_dual_cone}                                                                               \\
        \mathcal F_p
         & :=\{x\in K:T(x)=b\},
         & C_p                                                    & :=\inf_{x\in\mathcal F_p}\langle c|x\rangle,
        \label{eq:conic_primal}                                                                                  \\
        \mathcal F_d
         & :=\{y\in V':c-T^*(y)\in K^*\},
         & C_d                                                    & :=\sup_{y\in\mathcal F_d}\langle b|y\rangle.
        \label{eq:conic_dual}
    \end{align}
    The values belong to the extended real line.  The conventions are
    $\inf\varnothing=+\infty$ and $\sup\varnothing=-\infty$; an objective
    unbounded below has infimum $-\infty$, and one unbounded above has
    supremum $+\infty$.  The nonempty-interior assumption is Wolf's
    convention for a conic program; the definitions and weak duality do not
    require it.
\end{definition}

\begin{theorem}[Weak duality]
    \label{thm:conic_weak_duality}
    \lean{ConicProgram.weak_duality_pointwise, ConicProgram.weak_duality}
    \leanok
    \uses{def:conic_programs}
    For every $x\in\mathcal F_p$ and $y\in\mathcal F_d$,
    $\langle b|y\rangle\leq\langle c|x\rangle$.  Consequently,
    $C_d\leq C_p$ with the extended-real conventions of
    Definition~\ref{def:conic_programs}.
\end{theorem}

\begin{proof}
    \leanok
    \uses{def:conic_programs}
    Dual feasibility and $x\in K$ give
    \begin{align}
        0
         & \leq\langle x|c-T^*(y)\rangle
        =\langle c|x\rangle-\langle T(x)|y\rangle
        =\langle c|x\rangle-\langle b|y\rangle.
        \label{eq:conic_pointwise_weak_duality}
    \end{align}
    This is the pointwise inequality.  Taking the supremum over
    $y\in\mathcal F_d$ and then the infimum over $x\in\mathcal F_p$ proves
    $C_d\leq C_p$.
\end{proof}
